\section{Síntesis y ampliación}

\subsection{Un eje conductor transferible a otros sistemas generativos}

El método más transferible de esta clase no consiste en copiar los nombres de los modelos de Movie Gen, sino en diseñar el sistema siguiendo el orden de sus dependencias: primero se define la variable aleatoria que se va a generar y después se elige el objetivo de transporte; primero se calcula la longitud en token y después se eligen la attention y el paralelismo; primero se definen las dimensiones de evaluación y después se revisan los ejemplos y el scaling; por último, las aplicaciones se descomponen en distintos problemas condicionales: generación base, edición, identidad y audio.

\subsection{Tres preguntas de investigación abiertas}

Primera: ¿es posible comprimir aún más el latent espaciotemporal o usar attention dispersa o jerárquica sin sacrificar el movimiento de grano fino, para superar así el costo cuadrático de los videos largos? Segunda: ¿qué tipo de estado intermedio debería producir el reasoning sobre medios y cómo se construye un verifier coherente con el juicio humano? Tercera: ¿puede un modelo conjunto nativo de audio y video aprovechar la información entre modalidades y, al mismo tiempo, evitar la escasez de datos alineados y la interferencia entre tareas?

En conjunto, estas preguntas muestran que la siguiente generación de modelos de video no dependerá solo de agrandar el denoiser. Es más probable que los avances provengan del diseño coordinado de representation, data, attention, parallelism, verification y multimodal objective. Al evaluar un sistema nuevo, la pregunta más valiosa no es «qué tan espectaculares son los videos que genera», sino «en qué workload, con qué conjunto de evidencias y dentro de qué límite de costo modifica la Pareto frontier».

\subsection{Lecturas adicionales}

\begin{itemize}[itemsep=0.2em,topsep=0.2em]
  \item Grabación oficial: \url{https://www.youtube.com/watch?v=YGHF8_tf--g}
  \item Página del curso CS25 V5: \url{https://web.stanford.edu/class/cs25/past/cs25-v5/}
  \item Artículo original de Movie Gen: \url{https://ai.meta.com/static-resource/movie-gen-research-paper/}
  \item Temas para profundizar: Flow Matching, DiT/AdaLN, Temporal Autoencoder, context parallelism, human evaluation y media reasoning verification.
\end{itemize}

\end{document}
